$v_p(N)$で正の整数$N$が$p$で割れる回数を表すとする.

(1): $n=\sqrt[abc]{a!+b!+c!}$とおく. 明らかに$n>1$なので$n$には素因数$p$が存在する. いま
\[
v_p(c!) = \sum_{k=1}^{\infty} \parenc{\dfrac{c}{p^k}} < \sum_{k=1}^{\infty} \dfrac{c}{2^k} = c
\]
に注意すると, $n^{abc} = a!+b!+c!$より
\[v_p(a! + b! + c!) = v_p(n^{abc}) = abc v_p(n) > c > v_p(c!)
\]
だから, $a! + b! + c!$ は $p^{v_p(c!) + 1}$で割り切れる. これと,  $c!$は$p^{v_p(c!) + 1}$で割り切れないが$p^{v_p(c!)}$で割り切れることから, $a! + b! = (a! + b! + c!) - c!$もそうである. よって割り切れる回数は等しい. \par
$c\geq 3$とする. このとき, $n^{abc} = a! + b! + c! \leq 3\cdot c! = 3\cdot 2\cdot 3\cdot 4\cdot \cdots \dot c < c^{c}$である. よって$n^{ab} < c$, とくに,\bolm{p^{ab} < c}を得る. \\
(2): まず$b\geq 2$とする. $c=2$のときは$b=2$で, $n^{abc}$は$1! + 2! + 2! = 5$または$2!+2!+2! = 6$となるが, このとき$n$は正の整数とならないので不適. 以降$c\geq 3$とすると, (1)から
\[
v_p(c!) \geq v_p((p^{ab})!) \geq \dfrac{p^{ab}}{p} + \dfrac{p^{ab}}{p^{ab}} > p^{ab-1}
\]
を得る($ab\geq 2$に注意). これと(1)より$a! + b!$を評価できる. すなわち,
\[
a! + b! \geq p^{v_p(a!+b!)} = p^{v_p(c!)} > p^{(p^{ab-1})} \geq 2^{(2^{ab-1})}
\]
とくに$1\leq a\leq b$から
\[
2\cdot b! > 2^{(2^{b-1})} \implies \dfrac{b!}{2^{2^{b-1} - 1}} > 1
\]
が必要である. ここで$b\geq 2$において $2^{b-1} - 1 = 2^0 + 2^1 + \dots + 2^{b-2}$ なので
\[
\dfrac{b!}{2^{(2^{b-2}-1)}} = \dfrac{b}{2^{2^{b-2}}} \cdot \dfrac{b-1}{2^{2^{b-3}}} \cdot \cdots \cdot \dfrac{2}{2^{2^0}}
\]
とも表せ, 各因子の$\dfrac{k+2}{2^{2^{k}}}$ ($k\geq 0$)は $2^{2^k} \geq 2^{k+1} \geq k+2$ (二項定理による評価)より$1$以下. よって上の値は$1$以上になりえず, 不適.

以上より$b=1$が必要である. このとき$a=1$で, $n^c = 2 + c!$となる. $c\geq 4$とすると$n$は偶数であって, $n^c$は$4$の倍数だから$c! = n^c - 2$は4で割れない奇数である. これは不適. よって$c\leq 3$であり, $c=1,2,3$のとき$n=\sqrt[c]{2+c!} = 3,2,2$なので十分. 以上より, 求める組は
\[
\bolm{(a,b,c) = (1,1,1), (1,1,2), (1,1,3)}
\]
